% Chapter1_Introduction.tex-ээс 2026-11-27-нд үгчлэн зөөсөн агентын тэмдэглэл (хэвлэгдэхгүй). Шинэ тэмдэглэлийг энд нэмнэ.

%======================================================================
% ЭНЭ ФАЙЛД БАЙГАА БҮХ ЗҮЙЛ БАГШИД ҮЗҮҮЛЭХ НООРОГТ ХЭВЛЭГДЭНЭ (бүтэн ноорогт ч мөн).
% Зөвхөн бүтэн ноорогт орох хэсэг тусдаа файлд: Chapters/full/ch1_*.tex (\fullinput мөрүүд).
%
% 2026.11.16 (2), хэрэглэгчийн шүүмжээр дахин бичсэн: үндэслэлийг вэб сайтаар биш, сургуульд
% ирсэн хүн ба ажилтны бэрхшээлээр эхлүүлсэн; авсан мэдээллийг тоочихгүй (тоо 4.1-ийн хүснэгтэд);
% мэдээллийн эх сурвалж = вэб сайт ба сургуулийн албад; «мэдлэгийн мод» нь арга тул зорилт,
% ач холбогдлоос авсан (3.3-т тайлбарласан). Өмнөх хувилбар: thesis/archive-20261116/.
% Бүтэц, үг хэллэг: papers/28-mn-thesis-style-20261111.mn.md (Монгол 7 дипломын тулгалт).
%
% 2026.11.20 (хэрэглэгч: «Монгол судалгааны ажлын үг хэллэгийг судлаад бичээд үз»): 1.1, 1.3-ын найруулгыг
% Монгол эрдэм шинжилгээний бичвэрийн хүрээгээр зассан туршилтын хувилбар. Баримт, тоо, эшлэл, хүснэгт, шошго
% өөрчлөгдөөгүй; 1.2 (зорилго, зорилт)-д гар хүрээгүй. Өмнө/дараа: reports/Хэллэгийн туршилт 2026-11-20.md.
% Өмнөх хувилбар: thesis/archive-20261120/.
%
% 2026.11.23 (хэрэглэгч: «thesis-ийг эхнээс дуустал дахин харж, үг сонголтыг зас»): аудитаар
% (papers/45-thesis-language-audit-20261122.mn.md №10–32) засав — ШМТК-ийн товчлол, үндэслэлийн
% 2-р өгүүлбэрийг сайтаар хүрээлэхгүй, «хамгийн их» → «түгээмэл» (давтамж тоолоогүй), Facebook-ийн
% догол мөрөнд төлбөрийн зөрүү ба «нөхнө», \acs{WER}, «тулгуурлан» давталт, «сургуулийн зүгээс өгсөн».
% Зорилго, зорилтын утга, гарчиг өөрчлөгдөөгүй. Өмнөх хувилбар: thesis/archive-20261123/.
% 2026.11.23 (2), орхигдсон аудитын санал: «Бурч … (International Burch University)» (№18, Милано-Бикокка-тай нэг
% хэв), «(hallucination)» нэр томьёо (№21).
%======================================================================

% 2026.11.16 (8): Орхонтунгалаг багшийн асуулт («Facebook чатботоос юугаараа ялгаатай вэ, заавал
% сургууль дээр ирж асуух хэрэг байна уу») — одоо байгаа сувагтай харьцах байр суурийг энд хэлсэн.
% Дэлгэрэнгүй: 2.4, 2.5; хэлэлцүүлэг: papers/37-distinctiveness-vs-facebook-chat-20261116.mn.md.

% Зорилт 1: карьерын төв, ресепшний ажилтантай уулзсаны дараа уулзсан албадыг нэрээр нь бич.
% 2026.11.16 (3), хэрэглэгч: мэдлэгийн мод бол RAG-ийн суурь — 2-р зорилт, ач холбогдолд буцааж оруулсан;
% үнэлгээ тусдаа 4-р зорилт хэвээр (үнэлгээг «мэдлэгийн модонд тулгуурласан» гэж нэрлэхгүй).

% 2026.11.23: бүтэн ноорогт зураг бүлгийн төгсгөлд дангаараа хуудас эзэлж байсан тул зургийг тайлбарлах
% догол мөрийн өмнө байрлуулав (хоёр ноорогт хоёуланд нь бичвэрийн дунд орно).
% 2026.11.16 (багшийн зөвлөгөө): хэрэглээний тохиолдлын диаграммыг 3.2-оос энд, эхний бүлэгт шилжүүлсэн.
